\documentclass[11pt,a4paper]{article}

\usepackage[brazil]{babel}
\usepackage[utf8]{inputenc}
\usepackage[T1]{fontenc}
\usepackage{lmodern}
\usepackage{geometry}
\usepackage{setspace}
\usepackage{parskip}
\usepackage{titlesec}
\usepackage{xcolor}
\usepackage{hyperref}

\geometry{margin=2.5cm}
\setstretch{1.2}

\definecolor{primary}{RGB}{41,128,185}
\hypersetup{
    colorlinks=true,
    linkcolor=primary,
    urlcolor=primary
}

\titleformat{\section}{\Large\bfseries\color{primary}}{\thesection}{1em}{}

\title{\textbf{Declaração de Conhecimento \\ GULECD/UFPR}}
\author{Chapa 01}
\date{\today}

\begin{document}

\maketitle

\section{Declaração}

A presente chapa, identificada como \textbf{Chapa 01}, candidata à coordenação do Grupo de Usuários Linux e Software Livre do curso de Estatística e Ciência de Dados da Universidade Federal do Paraná (GULECD/UFPR), declara formalmente que:

\begin{itemize}
    \item possui pleno conhecimento do conteúdo da \textbf{Carta de Fundação} do GULECD/UFPR;
    \item compreende e está integralmente alinhada aos princípios estabelecidos no \textbf{Manifesto} do grupo;
    \item reconhece e compromete-se a cumprir as diretrizes do \textbf{Código de Conduta};
    \item tem ciência das normas estabelecidas no \textbf{Regulamento Eleitoral};
    \item está ciente das disposições da \textbf{Política de Privacidade};
    \item reconhece os demais documentos institucionais do grupo e compromete-se a respeitá-los integralmente.
\end{itemize}

A chapa declara ainda que se compromete a:

\begin{itemize}
    \item conduzir suas atividades em conformidade com os princípios de transparência, colaboração e responsabilidade coletiva;
    \item respeitar a cultura de software livre e os valores que orientam o GULECD/UFPR;
    \item atuar de forma ética, respeitosa e alinhada aos interesses da comunidade do grupo.
\end{itemize}

\section{Responsabilidade}

A presente declaração expressa o compromisso formal de todos os membros da chapa com os princípios institucionais do GULECD/UFPR, sendo parte integrante do processo eleitoral e da validação da candidatura.

\vspace{2cm}

\noindent\textbf{Coordenador Representativo:} \\
Nome: Marcos Oliveira de Carvalho \\


\vspace{1.5cm}

\noindent\textbf{Vice-Coordenador Representativo:} \\
Nome: [Inserir Nome Completo] \\


\vspace{2cm}

\noindent\textbf{Demais Membros da Chapa:} \\{}
\noindent[Inserir lista completa dos membros ou indicar documento complementar]

\end{document}